\documentclass[12pt,a4paper]{article}
\usepackage[utf8]{inputenc}
\usepackage[french]{babel}
\usepackage[T1]{fontenc}
\usepackage{geometry}
\geometry{margin=2.5cm, top=3cm}
\usepackage{xcolor}
\usepackage{titlesec}
\usepackage{listings}
\usepackage{tikz}
\usepackage{graphicx}
\usepackage[most]{tcolorbox}
\tcbuselibrary{listings, skins, breakable}

% Clean monospace font for code
\usepackage{inconsolata}
\usepackage{textcomp}

% ---- Colors ----
\definecolor{MainBlue}{RGB}{20, 60, 110}
\definecolor{LightBlue}{RGB}{235, 245, 255}
\definecolor{TerminalBlack}{RGB}{40, 44, 52}
\definecolor{ShadowGray}{RGB}{210, 210, 210}
\definecolor{CodeKeyword}{RGB}{86, 156, 214}
\definecolor{CodeString}{RGB}{206, 145, 80}
\definecolor{CodeComment}{RGB}{106, 153, 85}

% ---- Python listing style ----
\lstdefinestyle{pystyle}{
    language=Python,
    basicstyle=\ttfamily\footnotesize,
    keywordstyle=\color{CodeKeyword}\bfseries,
    stringstyle=\color{CodeString},
    commentstyle=\color{CodeComment}\itshape,
    numberstyle=\ttfamily\tiny\color{gray!60},
    columns=fixed,
    keepspaces=true,
    upquote=true,
    showstringspaces=false,
    numbers=left,
    numbersep=10pt,
    tabsize=4,
    breaklines=true,
    breakatwhitespace=false,
    breakindent=2em,
    postbreak=\mbox{\textcolor{gray}{\ensuremath{\hookrightarrow}}\space},
}

% ---- Section Title Format ----
\titleformat{\section}
  {\color{MainBlue}\normalfont\Large\bfseries}{}{0em}
  {%
    \begin{tikzpicture}[baseline=(char.base)]
      \fill[MainBlue] (0,-0.1) rectangle (0.3,0.6);
      \node (char) at (0.5,0) {};
    \end{tikzpicture}\hspace{0pt}%
  }

% ---- Code Box (uses newtcblisting for proper listing support) ----
\newtcblisting{codeholder}[1]{
    listing only,
    listing style=pystyle,
    enhanced,
    breakable,
    colback=white,
    colframe=TerminalBlack,
    boxrule=0.7pt,
    arc=4pt,
    left=2pt,
    right=4pt,
    top=4pt,
    bottom=4pt,
    boxsep=4pt,
    title={%
        \textcolor{red!80}{$\bullet$}
        \textcolor{yellow!80}{$\bullet$}
        \textcolor{green!80}{$\bullet$}
        \hspace{0.4cm}
        {\ttfamily\footnotesize #1}
    },
    fonttitle=\bfseries\color{white},
    colbacktitle=TerminalBlack,
    attach boxed title to top left={yshift=-3mm, xshift=6mm},
    boxed title style={size=small, sharp corners, arc=2pt},
    shadow={1mm}{-1mm}{0mm}{ShadowGray},
}

% ---- Exercise Command ----
\newcommand{\exercice}[1]{
    \clearpage
    \section{Exercice : #1}
}

\begin{document}

% ---- Cover Page ----
\begin{titlepage}
    \begin{tikzpicture}[remember picture, overlay]
        \fill[MainBlue]
            (current page.north west)
            rectangle ([xshift=1.5cm]current page.south west);
        \node[
            anchor=north west,
            fill=LightBlue,
            minimum width=15cm,
            minimum height=4cm,
            inner sep=1cm
        ] at ([xshift=2cm, yshift=-4cm]current page.north west) {
            \begin{minipage}{14cm}
                {\Huge\bfseries\color{MainBlue}
                    RAPPORT DE TRAVAUX PRATIQUES}\\[0.4cm]
                {\large\color{black!70}
                    Série d'exercices --- Perceptron / Deep Learning}
            \end{minipage}
        };
    \end{tikzpicture}

    \vfill

    \begin{flushright}
    \begin{tcolorbox}[
        width=8.5cm,
        colframe=MainBlue, colback=white,
        boxrule=2pt, arc=10pt, inner sep=15pt,
        shadow={3mm}{-3mm}{0mm}{ShadowGray}
    ]
        {\large\textbf{IDENTITÉ}}\vspace{5pt}\hrule\vspace{10pt}
        \begin{tabular}{ll}
            \textbf{Nom :}        & Jalal Grini \\
            \textbf{Classe :}     & TDIA2 \\
            \textbf{École :}      & ENSAH \\
            \textbf{Enseignant :} & [Nom du Professeur] \\
            \textbf{Sujet :}      & Perceptron -- Série d'exercices
        \end{tabular}
    \end{tcolorbox}
    \end{flushright}
    \vspace{2cm}
\end{titlepage}

\newpage
\tableofcontents

% ======================== EXERCICE 1 ========================
\exercice{Classification des portes logiques AND et OR}

\subsection*{Énoncé}
Apprendre au Perceptron les fonctions logiques AND et OR en se
conformant aux étapes suivantes :
\begin{enumerate}
    \item Créez les ensembles de données pour les portes logiques AND et OR.
    \item Entraînez un modèle Perceptron sur les données pour chaque porte séparément.
    \item Testez le modèle sur des données de validation.
    \item Commentez la capacité du Perceptron à apprendre les deux types de portes.
\end{enumerate}

\subsection*{1) Génération des données}

\begin{codeholder}{ex1.py}
# Données communes à toutes les implémentations
X = np.array([[0, 0], [0, 1], [1, 0], [1, 1]])
y_or = np.array([0, 1, 1, 1])
y_and = np.array([0, 0, 0, 1])
\end{codeholder}

\paragraph{Explication}
Le bloc ci-dessus crée le jeu de données d'entraînement pour les portes
logiques (toutes les combinaisons binaires de deux bits).
\texttt{X} contient les entrées ; \texttt{y\_or} et \texttt{y\_and}
sont les labels pour OR et AND respectivement.

% --

\subsection*{2) Fonction d'activation}

\begin{codeholder}{ex1.py}
def fct_activattion(c):
    """Fonction d'activation seuil (step function)"""
    if c >= 0.1:
        return 1
    else:
        return 0

def step_activation(x):
    """Step function for perceptron (vectorized)"""
    return (x >= 0).astype(int)
\end{codeholder}

\paragraph{Explication}
La fonction d'activation est un seuil (step function) fixé à 0.1.
Elle renvoie 1 si la somme pondérée est supérieure ou égale au seuil, sinon 0.
Nous utilisons également une version vectorisée de cette fonction pour l'implémentation Keras.

% --

\subsection*{3) Entraînement (from scratch)}

\begin{codeholder}{ex1.py}
def training(x, y, learn_rate=0.1, epochs=5):
    n = x.shape[1]
    w = np.zeros(n)
    b = 0
    for epoch in range(epochs):
        for i in range(len(x)):
            c = np.dot(x[i], w) + b
            y_pre = fct_activattion(c)
            err = y[i] - y_pre
            w += learn_rate * err * x[i]
            b += learn_rate * err
    return w, b

def prediction(x, w, b):
    prediction = []
    for i in range(len(x)):
        z = np.dot(x[i], w) + b
        prediction.append(fct_activattion(z))
    return np.array(prediction)
\end{codeholder}

\paragraph{Explication}
L'algorithme met en œuvre la règle du Perceptron : on initialise les poids
à zéro, puis pour chaque exemple et chaque époque on calcule l'erreur
et on met à jour les poids et le biais.

% --

\subsection*{4) Prédiction et évaluation (from scratch)}

\begin{codeholder}{ex1.py}
print("test for AND:\n")
w_and, b_and = training(X, y_and)
and_pre = prediction(X, w_and, b_and)
print("the weight is :", w_and, " and the bias is :", b_and)
print("the AND prediction is : ", and_pre)
print("the real AND is : ", y_and)

print("\ntest for OR:\n")
w_or, b_or = training(X, y_or)
or_pre = prediction(X, w_or, b_or)
print("the weight is :", w_or, " and the bias is :", b_or)
print("the OR prediction is : ", or_pre)
print("the real OR is : ", y_or)
\end{codeholder}

\paragraph{Explication}
La fonction \texttt{prediction} applique la combinaison linéaire puis
la fonction d'activation pour retourner les prédictions.
Le script teste ensuite AND puis OR (entraînement puis prédiction sur le même jeu).

% --

\subsection*{5) Implémentation avec scikit-learn}

\begin{codeholder}{ex1.py}
def sk_train(X, y):
    clf = perceptron.Perceptron(eta0=0.1, max_iter=10, shuffle=False)
    clf.fit(X, y)
    prediction = clf.predict(X)
    return clf, prediction
\end{codeholder}

\paragraph{Explication}
La classe \texttt{Perceptron} de \texttt{sklearn.linear\_model} encapsule
l'entraînement et fournit les poids via \texttt{clf.coef\_}
et le biais via \texttt{clf.intercept\_}.

% --

\subsection*{6) Résultats et comparaison avec scikit-learn}

\begin{codeholder}{ex1.py}
# Comparison between from-scratch and sklearn on TRAINING data
print("\n*** COMPARISON: From Scratch vs Sklearn (on TRAINING data) ***")
print("From scratch AND predictions:", and_pre)
print("Sklearn AND predictions:     ", and_sk_pre)
print("Match:", np.array_equal(and_pre, and_sk_pre))
print()
print("From scratch OR predictions: ", or_pre)
print("Sklearn OR predictions:      ", or_sk_pre)
print("Match:", np.array_equal(or_pre, or_sk_pre))

# VALIDATION ON SEPARATE DATA
print("\n" + "="*60)
print("VALIDATION ON NEW/TEST DATA")
print("="*60)

# Using the same data for validation (in this case, XOR-like validation to test limits)
# Same as training (complete dataset)
X_test = np.array([[0, 0], [0, 1], [1, 0], [1, 1]])

print("\n--- VALIDATION: And Gate ---")
and_pre_test = prediction(X_test, w_and, b_and)
print("AND Training Predictions: ", and_pre)
print("AND Validation Predictions:", and_pre_test)
print("Accuracy (Training): {:.2%}".format(np.mean(and_pre == y_and)))
print("Accuracy (Validation): {:.2%}".format(np.mean(and_pre_test == y_and)))

print("\n--- VALIDATION: OR Gate ---")
or_pre_test = prediction(X_test, w_or, b_or)
print("OR Training Predictions: ", or_pre)
print("OR Validation Predictions:", or_pre_test)
print("Accuracy (Training): {:.2%}".format(np.mean(or_pre == y_or)))
print("Accuracy (Validation): {:.2%}".format(np.mean(or_pre_test == y_or)))
\end{codeholder}

\paragraph{Explication}
Nous comparons maintenant les résultats de l'implémentation from scratch avec ceux de scikit-learn sur les données d'entraînement.
Nous validons également nos modèles sur un jeu de données séparé (dans ce cas, nous utilisons le même jeu de données pour simplifier l'exemple).
Les deux implémentations donnent des résultats identiques, ce qui confirme la correction de notre implémentation from scratch.

% --

\subsection*{2) Fonction d'activation}

\begin{codeholder}{ex1.py}
def fct_activattion(c):
    """Fonction d'activation seuil (step function)"""
    if c >= 0.1:
        return 1
    else:
        return 0

def step_activation(x):
    """Step function for perceptron (vectorized)"""
    return (x >= 0).astype(int)
\end{codeholder}

\paragraph{Explication}
La fonction d'activation est un seuil (step function) fixé à 0.1.
Elle renvoie 1 si la somme pondérée est supérieure ou égale au seuil, sinon 0.
Nous utilisons également une version vectorisée de cette fonction pour l'implémentation Keras.

% --

\subsection*{3) Entraînement (from scratch)}

\begin{codeholder}{ex1.py}
def training(x, y, learn_rate=0.1, epochs=5):
    n = x.shape[1]
    w = np.zeros(n)
    b = 0
    for epoch in range(epochs):
        for i in range(len(x)):
            c = np.dot(x[i], w) + b
            y_pre = fct_activattion(c)
            err = y[i] - y_pre
            w += learn_rate * err * x[i]
            b += learn_rate * err
    return w, b

def prediction(x, w, b):
    prediction = []
    for i in range(len(x)):
        z = np.dot(x[i], w) + b
        prediction.append(fct_activattion(z))
    return np.array(prediction)
\end{codeholder}

\paragraph{Explication}
L'algorithme met en œuvre la règle du Perceptron : on initialise les poids
à zéro, puis pour chaque exemple et chaque époque on calcule l'erreur
et on met à jour les poids et le biais.

% --

\subsection*{4) Prédiction et évaluation (from scratch)}

\begin{codeholder}{ex1.py}
print("test for AND:\n")
w_and, b_and = training(X, y_and)
and_pre = prediction(X, w_and, b_and)
print("the weight is :", w_and, " and the bias is :", b_and)
print("the AND prediction is : ", and_pre)
print("the real AND is : ", y_and)

print("\ntest for OR:\n")
w_or, b_or = training(X, y_or)
or_pre = prediction(X, w_or, b_or)
print("the weight is :", w_or, " and the bias is :", b_or)
print("the OR prediction is : ", or_pre)
print("the real OR is : ", y_or)
\end{codeholder}

\paragraph{Explication}
La fonction \texttt{prediction} applique la combinaison linéaire puis
la fonction d'activation pour retourner les prédictions.
Le script teste ensuite AND puis OR (entraînement puis prédiction sur le même jeu).

% --

\subsection*{5) Implémentation avec scikit-learn}

\begin{codeholder}{ex1.py}
def sk_train(X, y):
    clf = perceptron.Perceptron(eta0=0.1, max_iter=10, shuffle=False)
    clf.fit(X, y)
    prediction = clf.predict(X)
    return clf, prediction
\end{codeholder}

\paragraph{Explication}
La classe \texttt{Perceptron} de \texttt{sklearn.linear\_model} encapsule
l'entraînement et fournit les poids via \texttt{clf.coef\_}
et le biais via \texttt{clf.intercept\_}.

% --

\subsection*{6) Résultats et comparaison avec scikit-learn}

\begin{codeholder}{ex1.py}
# Comparison between from-scratch and sklearn on TRAINING data
print("\n*** COMPARISON: From Scratch vs Sklearn (on TRAINING data) ***")
print("From scratch AND predictions:", and_pre)
print("Sklearn AND predictions:     ", and_sk_pre)
print("Match:", np.array_equal(and_pre, and_sk_pre))
print()
print("From scratch OR predictions: ", or_pre)
print("Sklearn OR predictions:      ", or_sk_pre)
print("Match:", np.array_equal(or_pre, or_sk_pre))

# VALIDATION ON SEPARATE DATA
print("\n" + "="*60)
print("VALIDATION ON NEW/TEST DATA")
print("="*60)

# Using the same data for validation (in this case, XOR-like validation to test limits)
# Same as training (complete dataset)
X_test = np.array([[0, 0], [0, 1], [1, 0], [1, 1]])

print("\n--- VALIDATION: AND Gate ---")
and_pre_test = prediction(X_test, w_and, b_and)
print("AND Training Predictions: ", and_pre)
print("AND Validation Predictions:", and_pre_test)
print("Accuracy (Training): {:.2%}".format(np.mean(and_pre == y_and)))
print("Accuracy (Validation): {:.2%}".format(np.mean(and_pre_test == y_and)))

print("\n--- VALIDATION: OR Gate ---")
or_pre_test = prediction(X_test, w_or, b_or)
print("OR Training Predictions: ", or_pre)
print("OR Validation Predictions:", or_pre_test)
print("Accuracy (Training): {:.2%}".format(np.mean(or_pre == y_or)))
print("Accuracy (Validation): {:.2%}".format(np.mean(or_pre_test == y_or)))
\end{codeholder}

\paragraph{Explication}
Nous comparons maintenant les résultats de l'implémentation from scratch avec ceux de scikit-learn sur les données d'entraînement.
Nous validons également nos modèles sur un jeu de données séparé (dans ce cas, nous utilisons le même jeu de données pour simplifier l'exemple).
Les deux implémentations donnent des résultats identiques, ce qui confirme la correction de notre implémentation from scratch.

% --

\subsection*{7) Implémentation avec Keras}

\begin{codeholder}{ex1.py}
# Create a simple perceptron-like model using Keras
def create_perceptron_model():
    model = Sequential([
        Dense(1, input_shape=(2,), activation='linear', 
              kernel_initializer='zeros', bias_initializer='zeros')
    ])
    return model

# Train using perceptron learning rule (manual implementation with Keras model)
def train_perceptron(model, X, y, learning_rate=0.1, epochs=5):
    # Get weights and bias
    weights = model.get_weights()[0]  # weight matrix (2, 1)
    bias = model.get_weights()[1]     # bias vector (1,)
    
    # Perceptron learning rule
    for epoch in range(epochs):
        for i in range(len(X)):
            # Linear combination
            linear_output = np.dot(X[i], weights[:, 0]) + bias[0]  # Extract column 0
            # Step activation
            y_pred = step_activation(linear_output)
            # Error
            error = y[i] - y_pred
            # Update if error exists
            if error != 0:
                weights[:, 0] += learning_rate * error * X[i]  # Update column 0
                bias[0] += learning_rate * error
    
    # Update model weights
    model.set_weights([weights, bias])
    return model

def predict_perceptron(model, X):
    weights = model.get_weights()[0]  # (2, 1)
    bias = model.get_weights()[1]     # (1,)
    linear_output = np.dot(X, weights[:, 0]) + bias[0]  # Extract column 0
    return step_activation(linear_output)
\end{codeholder}

\paragraph{Explication}
Nous définissons ensuite notre modèle perceptron en utilisant Keras. Le modèle consiste en une seule couche dense avec 2 entrées (pour les deux caractéristiques) et 1 sortie.
Nous utilisons une activation linéaire car nous allons appliquer manuellement la fonction d'activation de seuil (step function) lors de l'entraînement,
ce qui nous permet de suivre exactement la même règle d'apprentissage que dans l'implémentation from scratch.

% --

\begin{codeholder}{ex1.py}
# Test AND gate
print("test for AND:\n")
model_and = create_perceptron_model()
model_and = train_perceptron(model_and, X, y_and, learning_rate=0.1, epochs=5)
and_pred = predict_perceptron(model_and, X)
weights_and = model_and.get_weights()[0][:, 0]  # Extract column 0 and flatten
bias_and = model_and.get_weights()[1][0]
print("the weight is :", weights_and, " and the bias is :", bias_and)
print("the AND prediction is : ", and_pred)
print("the real AND is : ", y_and)

print("\n-----------------------------------------")
print("test for OR:\n")
model_or = create_perceptron_model()
model_or = train_perceptron(model_or, X, y_or, learning_rate=0.1, epochs=5)
or_pred = predict_perceptron(model_or, X)
weights_or = model_or.get_weights()[0][:, 0]  # Extract column 0 and flatten
bias_or = model_or.get_weights()[1][0]
print("the weight is :", weights_or, " and the bias is :", bias_or)
print("the OR prediction is : ", or_pred)
print("the real OR is : ", y_or)

print("\n-----------------------------------------")
print("VALIDATION ON SAME DATA (for completeness)")
print("AND Validation Predictions:", predict_perceptron(model_and, X))
print("OR Validation Predictions:", predict_perceptron(model_or, X))
print("Accuracy (AND): {:.2%}".format(np.mean(predict_perceptron(model_and, X) == y_and)))
print("Accuracy (OR): {:.2%}".format(np.mean(predict_perceptron(model_or, X) == y_or)))
\end{codeholder}

\paragraph{Explication}
Finalement, nous testons notre modèle sur les portes AND et OR, tout comme dans les implémentations précédentes.
Nous affichons les poids et biais appris, ainsi que les prédictions du modèle comparées aux valeurs réelles.
Nous calculons également la précision sur les mêmes données d'entraînement (validation) pour compléter l'analyse.

% --

% --

\subsection*{6) Implémentation avec Keras}

\begin{codeholder}{ex1.py}
print("------------------keras method----------------")

X = np.array([[0, 0], [0, 1], [1, 0], [1, 1]])
y_or = np.array([0, 1, 1, 1])
y_and = np.array([0, 0, 0, 1])
\end{codeholder}

\paragraph{Explication}
Tout comme dans les implémentations précédentes, nous commençons par créer les ensembles de données pour les portes logiques AND et OR.
\texttt{X} contient les entrées ; \texttt{y\_or} et \texttt{y\_and} sont les labels pour OR et AND respectivement.

% --

\begin{codeholder}{ex1.py}
def step_activation(x):
    """Step function for perceptron"""
    return (x >= 0).astype(int)

# Create a simple perceptron-like model using Keras
def create_perceptron_model():
    model = Sequential([
        Dense(1, input_dim=2, activation='linear', 
              kernel_initializer='zeros', bias_initializer='zeros')
    ])
    return model
\end{codeholder}

\paragraph{Explication}
Nous définons ensuite notre modèle perceptron en utilisant Keras. Le modèle consiste en une seule couche dense avec 2 entrées (pour les deux caractéristiques) et 1 sortie.
Nous utilisons une activation linéaire car nous allons appliquer manuellement la fonction d'activation de seuil (step function) lors de l'entraînement,
ce qui nous permet de suivre exactement la même règle d'apprentissage que dans l'implémentation from scratch.

% --

\begin{codeholder}{ex1.py}
# Train using perceptron learning rule (manual implementation with Keras model)
def train_perceptron(model, X, y, learning_rate=0.1, epochs=5):
    # Get weights and bias
    weights = model.get_weights()[0]  # weight matrix (2, 1)
    bias = model.get_weights()[1]     # bias vector (1,)
    
    # Perceptron learning rule
    for epoch in range(epochs):
        for i in range(len(X)):
            # Linear combination
            linear_output = np.dot(X[i], weights[:, 0]) + bias[0]  # Extract column 0
            # Step activation
            y_pred = step_activation(linear_output)
            # Error
            error = y[i] - y_pred
            # Update if error exists
            if error != 0:
                weights[:, 0] += learning_rate * error * X[i]  # Update column 0
                bias[0] += learning_rate * error
    
    # Update model weights
    model.set_weights([weights, bias])
    return model

def predict_perceptron(model, X):
    weights = model.get_weights()[0]  # (2, 1)
    bias = model.get_weights()[1]     # (1,)
    linear_output = np.dot(X, weights[:, 0]) + bias[0]  # Extract column 0
    return step_activation(linear_output)
\end{codeholder}

\paragraph{Explication}
Nous implémentons ensuite la règle d'apprentissage du perceptron de manière manuelle, tout comme dans l'implémentation from scratch.
Cette fonction calcule la sortie linéaire, applique la fonction d'activation de seuil, calcule l'erreur,
et met à jour les poids et le biais selon la règle du perceptron.
Cette approche nous permet de comparer directement les trois méthodes en utilisant exactement le même algorithme d'apprentissage.

% --

\begin{codeholder}{ex1.py}
# Test AND gate
print("test for AND:\n")
model_and = create_perceptron_model()
model_and = train_perceptron(model_and, X, y_and, learning_rate=0.1, epochs=5)
and_pred = predict_perceptron(model_and, X)
weights_and = model_and.get_weights()[0][:, 0]  # Extract column 0 and flatten
bias_and = model_and.get_weights()[1][0]
print("the weight is :", weights_and, " and the bias is :", bias_and)
print("the AND prediction is : ", and_pred)
print("the real AND is : ", y_and)

print("\n-----------------------------------------")
print("test for OR:\n")
model_or = create_perceptron_model()
model_or = train_perceptron(model_or, X, y_or, learning_rate=0.1, epochs=5)
or_pred = predict_perceptron(model_or, X)
weights_or = model_or.get_weights()[0][:, 0]  # Extract column 0 and flatten
bias_or = model_or.get_weights()[1][0]
print("the weight is :", weights_or, " and the bias is :", bias_or)
print("the OR prediction is : ", or_pred)
print("the real OR is : ", y_or)

print("\n-----------------------------------------")
print("VALIDATION ON SAME DATA (for completeness)")
print("AND Validation Predictions:", predict_perceptron(model_and, X))
print("OR Validation Predictions:", predict_perceptron(model_or, X))
print("Accuracy (AND): {:.2%}".format(np.mean(predict_perceptron(model_and, X) == y_and)))
print("Accuracy (OR): {:.2%}".format(np.mean(predict_perceptron(model_or, X) == y_or)))
\end{codeholder}

\paragraph{Explication}
Finalement, nous testons notre modèle sur les portes AND et OR, tout comme dans les implémentations précédentes.
Nous affichons les poids et biais appris, ainsi que les prédictions du modèle comparées aux valeurs réelles.
Nous calculons également la précision sur les mêmes données d'entraînement (validation) pour compléter l'analyse.

% --

\subsection*{8) Analyse et conclusions}

\begin{codeholder}{ex1.py}
# COMMENTARY ON PERCEPTRON LEARNING ABILITY
print("\n" + "="*60)
print("ANALYSIS: PERCEPTRON LEARNING ABILITY")
print("="*60)
print("""
FINDINGS:

1. LINEAR SEPARABILITY:
   - The Perceptron successfully learns BOTH AND and OR gates
   - These functions are LINEARLY SEPARABLE,
     meaning a single straight line can separate them
   - The Perceptron is designed for linearly separable problems

2. AND GATE:
   - Only one input combination (1,1) produces output 1
   - The hyperplane learned separates this clearly
   - The Perceptron learns this perfectly ✓

3. OR GATE:
   - Three input combinations (0,1), (1,0), (1,1) produce output 1
   - The hyperplane separates "at least one 1" from "both 0"
   - The Perceptron learns this perfectly ✓

4. LIMITATIONS OF PERCEPTRON:
   - The Perceptron CANNOT learn XOR (exclusive OR) gate
   - XOR is NOT linearly separable (requires multiple layers)
   - This was a fundamental limitation discovered in the 1960s
   - Multi-layer neural networks are needed for XOR

5. CONVERGENCE:
   - With proper learning rate and epochs, both gates converge
   - The from-scratch and sklearn implementations match ✓
   - Training on the complete dataset = validation on the same data
   - (In practice, you would use train/test split for generalization testing)

CONCLUSION:
The Perceptron is effective for linearly separable logical functions like AND and OR.
For non-linearly separable functions (like XOR), deeper architectures are required.
""")
\end{codeholder}

\subsection*{Résultats numériques}
Les résultats imprimés par le script (exécution locale) sont les suivants :
\begin{itemize}
    \item AND --- Accuracy (Training) : 100.00\% ; Accuracy (Validation) : 100.00\%.
    \item OR  --- Accuracy (Training) : 100.00\% ; Accuracy (Validation) : 100.00\%.
\end{itemize}

\paragraph{Interprétation}
Le Perceptron apprend parfaitement les fonctions linéairement séparables
(AND, OR). La faiblesse sur la validation pour AND indique que la frontière
apprise n'est pas généralisable sur ce petit exemple synthétique.


% ======================== EXERCICE 2 ========================
\exercice{Classification binaire sur les données Iris Setosa}

\subsection*{Énoncé}
Chargez le jeu de données Iris. Sélectionnez deux caractéristiques
(longueur et largeur des pétales). Entraînez le Perceptron pour distinguer
Iris Setosa des autres catégories.
Calculez la précision du modèle sur des données de test.

\subsection*{1) Chargement et sélection des caractéristiques}

\begin{codeholder}{ex2.py}
import numpy as np
import matplotlib.pyplot as plt
from sklearn import datasets
from sklearn.model_selection import train_test_split
from sklearn.preprocessing import StandardScaler
from sklearn.linear_model import Perceptron
from sklearn.metrics import (
    accuracy_score, confusion_matrix, classification_report
)

# Load Iris Dataset
iris = datasets.load_iris()
X    = iris.data
y    = iris.target

# Select Petal Length and Petal Width
X_selected = X[:, [2, 3]]

# Create binary labels: Setosa (1) vs Others (0)
y_binary = (y == 0).astype(int)
\end{codeholder}

\paragraph{Explication}
On charge le dataset Iris via \texttt{sklearn.datasets} puis on sélectionne
deux caractéristiques (index 2 et 3 : longueur et largeur des pétales).
Les labels sont transformés en problème binaire (Setosa / Non-Setosa).

% --

\subsection*{2) Prétraitement et entraînement}

\begin{codeholder}{ex2.py}
# Split data into training and test sets
X_train, X_test, y_train, y_test = train_test_split(
    X_selected, y_binary,
    test_size=0.2, random_state=42, stratify=y_binary
)

# Standardize features
scaler         = StandardScaler()
X_train_scaled = scaler.fit_transform(X_train)
X_test_scaled  = scaler.transform(X_test)

# Train Perceptron
clf = Perceptron(eta0=0.1, max_iter=100, random_state=42, tol=1e-3)
clf.fit(X_train_scaled, y_train)
\end{codeholder}

\paragraph{Explication}
La standardisation est importante pour la convergence du Perceptron.
On utilise \texttt{StandardScaler} puis on entraîne \texttt{Perceptron}
de sklearn avec un taux d'apprentissage et un nombre d'itérations fixés.

% --

\subsection*{3) Évaluation et visualisation}

\begin{codeholder}{ex2.py}
# Make predictions
y_pred_train = clf.predict(X_train_scaled)
y_pred_test  = clf.predict(X_test_scaled)

# Evaluate model
train_accuracy = accuracy_score(y_train, y_pred_train)
test_accuracy  = accuracy_score(y_test,  y_pred_test)

print(f"Training Accuracy: {train_accuracy:.2%}")
print(f"Test Accuracy:     {test_accuracy:.2%}")

print("\nConfusion Matrix (Test Data):")
cm = confusion_matrix(y_test, y_pred_test)
print(cm)

print("\nClassification Report:")
print(classification_report(
    y_test, y_pred_test,
    target_names=['Non-Setosa', 'Setosa']
))
\end{codeholder}

\paragraph{Explication}
Les métriques utilisées comprennent la précision et la matrice de confusion.
Le script sauvegarde également une figure (\texttt{iris\_classification.png})
pour visualiser la frontière de décision.

% --

\subsection*{4) Implémentation avec Keras}

\begin{codeholder}{ex2.py}
import numpy as np
import matplotlib.pyplot as plt
from sklearn import datasets
from sklearn.model_selection import train_test_split
from sklearn.preprocessing import StandardScaler
from sklearn.metrics import accuracy_score, confusion_matrix, classification_report
import tensorflow as tf
from tensorflow.keras.models import Sequential
from tensorflow.keras.layers import Dense
from tensorflow.keras.optimizers import SGD

# Load Iris Dataset
iris = datasets.load_iris()
X = iris.data
y = iris.target

# Select Petal Length and Petal Width
X_selected = X[:, [2, 3]]

# Create binary labels: Setosa (1) vs Others (0)
y_binary = (y == 0).astype(int)
\end{codeholder}

\paragraph{Explication}
Comme dans les implémentations précédentes, nous commençons par charger le dataset Iris via \texttt{sklearn.datasets} puis nous sélectionnons deux caractéristiques (index 2 et 3 : longueur et largeur des pétales).
Les labels sont transformés en problème binaire (Setosa / Non-Setosa).

% --

\begin{codeholder}{ex2.py}
# Split data into training and test sets
X_train, X_test, y_train, y_test = train_test_split(
    X_selected, y_binary, test_size=0.2, random_state=42, stratify=y_binary
)

# Standardize features
scaler = StandardScaler()
X_train_scaled = scaler.fit_transform(X_train)
X_test_scaled = scaler.transform(X_test)
\end{codeholder}

\paragraph{Explication}
La standardisation est importante pour la convergence du Perceptron.
Nous utilisons \texttt{StandardScaler} puis nous préparons les données pour l'entraînement et le test.

% --

\begin{codeholder}{ex2.py}
# Create and train Perceptron using Keras
def create_perceptron_model(input_dim):
    model = Sequential([
        Dense(1, input_dim=input_dim, activation='sigmoid',  # Using sigmoid for binary classification
              kernel_initializer='zeros', bias_initializer='zeros')
    ])
    optimizer = SGD(learning_rate=0.1)
    model.compile(optimizer=optimizer, loss='binary_crossentropy', metrics=['accuracy'])
    return model

# Train the model
model = create_perceptron_model(X_train_scaled.shape[1])
history = model.fit(X_train_scaled, y_train, 
                    epochs=100, 
                    verbose=0,
                    validation_split=0.1)
\end{codeholder}

\paragraph{Explication}
Nous définons ensuite notre modèle perceptron en utilisant Keras avec une activation sigmoïde adaptée à la classification binaire.
Nous utilisons l'optimiseur SGD avec un taux d'apprentissage de 0.1 et optimisons la perte binaire croisée.
Après avoir créé le modèle, nous l'entraînons sur les données d'entraînement standardisées.

% --

\begin{codeholder}{ex2.py}
# Make predictions
y_pred_train_prob = model.predict(X_train_scaled)
y_pred_train = (y_pred_train_prob > 0.5).astype(int).flatten()
y_pred_test_prob = model.predict(X_test_scaled)
y_pred_test = (y_pred_test_prob > 0.5).astype(int).flatten()

# Evaluate model
train_accuracy = accuracy_score(y_train, y_pred_train)
test_accuracy = accuracy_score(y_test, y_pred_test)

print(f"Training Accuracy: {train_accuracy:.2%}")
print(f"Test Accuracy: {test_accuracy:.2%}")

print("\nConfusion Matrix (Test Data):")
cm = confusion_matrix(y_test, y_pred_test)
print(cm)

print("\nClassification Report:")
print(classification_report(y_test, y_pred_test,
      target_names=['Non-Setosa', 'Setosa']))

# Get weights and bias (for comparison with sklearn)
weights = model.get_weights()[0].flatten()  # weights for the two features
bias = model.get_weights()[1][0]           # bias
print(f"\nLearned weights: {weights}")
print(f"Learned bias: {bias}")
\end{codeholder}

\paragraph{Explication}
Nous faisons ensuite des prédictions sur les ensembles d'entraînement et de test, puis nous évaluons le modèle en utilisant les mêmes métriques que précédemment : précision, matrice de confusion et rapport de classification.
Nous extrayons également les poids et le biais appris pour les comparer avec ceux obtenus par scikit-learn.

% --

\begin{codeholder}{ex2.py}
# Visualization (similar to original)
h = 0.02
x_min, x_max = X_test_scaled[:, 0].min() - 1, X_test_scaled[:, 0].max() + 1
y_min, y_max = X_test_scaled[:, 1].min() - 1, X_test_scaled[:, 1].max() + 1
xx, yy = np.meshgrid(np.arange(x_min, x_max, h), np.arange(y_min, y_max, h))

Z = model.predict(np.c_[xx.ravel(), yy.ravel()])
Z = (Z > 0.5).astype(int).reshape(xx.shape)

plt.figure(figsize=(12, 5))

# Training data plot
plt.subplot(1, 2, 1)
plt.contourf(xx, yy, Z, alpha=0.3, cmap='RdYlBu')
plt.scatter(X_train_scaled[y_train == 0, 0], X_train_scaled[y_train == 0, 1],
            c='red', marker='o', label='Non-Setosa', edgecolors='k')
plt.scatter(X_train_scaled[y_train == 1, 0], X_train_scaled[y_train == 1, 1],
            c='blue', marker='s', label='Setosa', edgecolors='k')
plt.xlabel('Petal Length')
plt.ylabel('Petal Width')
plt.title(f'Training Data (Accuracy: {train_accuracy:.1%})')
plt.legend()
plt.grid(True, alpha=0.3)

# Test data plot
plt.subplot(1, 2, 2)
plt.contourf(xx, yy, Z, alpha=0.3, cmap='RdYlBu')
correct = y_test == y_pred_test
plt.scatter(X_test_scaled[correct & (y_test == 0), 0],
            X_test_scaled[correct & (y_test == 0), 1],
            c='red', marker='o', label='Non-Setosa (Correct)', edgecolors='k')
plt.scatter(X_test_scaled[correct & (y_test == 1), 0],
            X_test_scaled[correct & (y_test == 1), 1],
            c='blue', marker='s', label='Setosa (Correct)', edgecolors='k')
incorrect = y_test != y_pred_test
plt.scatter(X_test_scaled[incorrect, 0], X_test_scaled[incorrect, 1],
            c='yellow', marker='x', s=200, label='Misclassified', linewidths=3)
plt.xlabel('Petal Length')
plt.ylabel('Petal Width')
plt.title(f'Test Data (Accuracy: {test_accuracy:.1%})')
plt.legend()
plt.grid(True, alpha=0.3)

plt.tight_layout()
plt.savefig('iris_classification.png', dpi=100, bbox_inches='tight')
plt.show()
\end{codeholder}

\paragraph{Explication}
Finalement, nous visualisons les résultats de la même manière que dans l'implémentation scikit-learn pour faciliter la comparaison.
Nous affichons la frontière de décision ainsi que la répartition des points d'entraînement et de test, avec une distinction entre les classifications correctes et incorrectes.

% --

\subsection*{Résultats et visualisation}

\begin{figure}[htbp]
    \centering
    \includegraphics[width=0.9\textwidth]{iris_classification.png}
    \caption{Frontière de décision et répartition des points
             (Iris Setosa vs autres).}
    \label{fig:iris_classif}
\end{figure}

\paragraph{Résultats chiffrés}
Lors de l'exécution : \textbf{Training Accuracy = 100.00\%},
\textbf{Test Accuracy = 100.00\%}.
La matrice de confusion sur le jeu de test est :
\begin{verbatim}
[[20  0]
 [ 0 10]]
\end{verbatim}

\paragraph{Interprétation}
La séparation est excellente et le Perceptron linéaire suffit pour
distinguer Setosa des autres classes lorsque l'on utilise uniquement
les deux caractéristiques des pétales.


% ======================== EXERCICE 3 ========================
\exercice{Régression linéaire avec le Perceptron}

\subsection*{Énoncé}
Créez un ensemble de données en 2 dimensions (x,y) où y=2x+1.
Divisez les données en ensembles d'entraînement et de test.
Utilisez le Perceptron pour prédire la valeur de y en fonction de x.
Évaluez le modèle en utilisant le Mean Squared Error.

\subsection*{1) Création des données}

\begin{codeholder}{ex3.py}
import numpy as np
import matplotlib.pyplot as plt
from sklearn.model_selection import train_test_split
from sklearn.linear_model import SGDRegressor
from sklearn.metrics import (
    mean_squared_error, mean_absolute_error, r2_score
)
from sklearn.preprocessing import StandardScaler

# Create dataset: y = 2x + 1 with some noise
np.random.seed(42)
X = np.random.uniform(-10, 10, 100).reshape(-1, 1)
y = 2 * X.flatten() + 1 + np.random.normal(0, 2, 100)

# Split data into training and test sets
X_train, X_test, y_train, y_test = train_test_split(
    X, y, test_size=0.2, random_state=42
)

# Standardize features
scaler         = StandardScaler()
X_train_scaled = scaler.fit_transform(X_train)
X_test_scaled  = scaler.transform(X_test)
\end{codeholder}

\paragraph{Explication}
Le jeu de données synthétique est créé selon la relation linéaire avec bruit.
On standardise ensuite les caractéristiques pour l'entraînement.

% --

\subsection*{3) Implémentation avec Keras}

\begin{codeholder}{ex3.py}
import numpy as np
import matplotlib.pyplot as plt
from sklearn.model_selection import train_test_split
from sklearn.preprocessing import StandardScaler
from sklearn.metrics import mean_squared_error, mean_absolute_error, r2_score
import tensorflow as tf
from tensorflow.keras.models import Sequential
from tensorflow.keras.layers import Dense
from tensorflow.keras.optimizers import SGD

# Create dataset: y = 2x + 1 with some noise
np.random.seed(42)
X = np.random.uniform(-10, 10, 100).reshape(-1, 1)
y = 2 * X.flatten() + 1 + np.random.normal(0, 2, 100)

# Split data into training and test sets
X_train, X_test, y_train, y_test = train_test_split(
    X, y, test_size=0.2, random_state=42
)

# Standardize features
scaler = StandardScaler()
X_train_scaled = scaler.fit_transform(X_train)
X_test_scaled = scaler.transform(X_test)
\end{codeholder}

\paragraph{Explication}
Comme dans les implémentations précédentes, nous commençons par créer le jeu de données synthétique selon la relation linéaire avec bruit.
Nous divisons ensuite les données en ensembles d'entraînement et de test, puis nous standardisons les caractéristiques pour l'entraînement.

% --

\begin{codeholder}{ex3.py}
# Create and train Perceptron using Keras (for regression, we use linear activation)
def create_perceptron_model(input_dim):
    model = Sequential([
        Dense(1, input_dim=input_dim, activation='linear',  # Linear activation for regression
              kernel_initializer='zeros', bias_initializer='zeros')
    ])
    optimizer = SGD(learning_rate=0.01)
    model.compile(optimizer=optimizer, loss='mean_squared_error', metrics=['mae'])
    return model

# Train the model
model = create_perceptron_model(X_train_scaled.shape[1])
history = model.fit(X_train_scaled, y_train, 
                    epochs=1000, 
                    verbose=0,
                    validation_split=0.1)
\end{codeholder}

\paragraph{Explication}
Nous définons ensuite notre modèle perceptron en utilisant Keras avec une activation linéaire adaptée à la régression.
Nous utilisons l'optimiseur SGD avec un taux d'apprentissage de 0.01 et optimisons l'erreur quadratique moyenne (MSE).
Après avoir créé le modèle, nous l'entraînons sur les données d'entraînement standardisées.

% --

\begin{codeholder}{ex3.py}
# Make predictions
y_pred_train = model.predict(X_train_scaled).flatten()
y_pred_test = model.predict(X_test_scaled).flatten()

# Evaluate model
mse_train = mean_squared_error(y_train, y_pred_train)
mse_test = mean_squared_error(y_test, y_pred_test)
mae_test = mean_absolute_error(y_test, y_pred_test)
r2_test = r2_score(y_test, y_pred_test)

print(f"Training MSE: {mse_train:.4f}")
print(f"Test MSE: {mse_test:.4f}")
print(f"Test MAE: {mae_test:.4f}")
print(f"R² Score: {r2_test:.4f}")

print(f"\nLearned weight: {model.get_weights()[0][0][0]:.4f}")
print(f"Learned bias: {model.get_weights()[1][0]:.4f}")
\end{codeholder}

\paragraph{Explication}
Nous faisons ensuite des prédictions sur les ensembles d'entraînement et de test, puis nous évaluons le modèle en utilisant les mêmes métriques que précédemment : MSE, MAE et R².
Nous extrayons également le poids et le biais appris pour les comparer avec ceux obtenus par SGDRegressor.

% --

\begin{codeholder}{ex3.py}
# Visualization
plt.figure(figsize=(12, 5))

# Training data plot
plt.subplot(1, 2, 1)
plt.scatter(X_train_scaled, y_train, c='blue',
            alpha=0.6, label='Training data')
plt.plot(X_train_scaled, y_pred_train, c='red',
         linewidth=2, label='Fitted line')
plt.xlabel('X (standardized)')
plt.ylabel('y')
plt.title(f'Training Data (MSE: {mse_train:.4f})')
plt.legend()
plt.grid(True, alpha=0.3)

# Test data plot
plt.subplot(1, 2, 2)
plt.scatter(X_test_scaled, y_test, c='blue', alpha=0.6, label='Test data')
plt.plot(X_test_scaled, y_pred_test, c='red', linewidth=2, label='Fitted line')
plt.xlabel('X (standardized)')
plt.ylabel('y')
plt.title(f'Test Data (MSE: {mse_test:.4f})')
plt.legend()
plt.grid(True, alpha=0.3)

plt.tight_layout()
plt.savefig('regression.png', dpi=100, bbox_inches='tight')
plt.show()
\end{codeholder}

\paragraph{Explication}
Finalement, nous visualisons les résultats de la même manière que dans l'implémentation scikit-learn pour faciliter la comparaison.
Nous affichons les données d'entraînement et de test avec la droite ajustée par le modèle, ainsi que les métriques d'erreur correspondantes.

% --

\subsection*{2) Entraînement et évaluation}

\begin{codeholder}{ex3.py}
# Train Perceptron (SGDRegressor for regression)
regressor = SGDRegressor(
    loss='squared_error',
    eta0=0.01,
    max_iter=1000,
    random_state=42
)
regressor.fit(X_train_scaled, y_train)

# Make predictions
y_pred_train = regressor.predict(X_train_scaled)
y_pred_test  = regressor.predict(X_test_scaled)

# Evaluate model
mse_train = mean_squared_error(y_train,  y_pred_train)
mse_test  = mean_squared_error(y_test,   y_pred_test)
mae_test  = mean_absolute_error(y_test,  y_pred_test)
r2_test   = r2_score(y_test,             y_pred_test)

print(f"Training MSE: {mse_train:.4f}")
print(f"Test MSE:     {mse_test:.4f}")
print(f"Test MAE:     {mae_test:.4f}")
print(f"R2 Score:     {r2_test:.4f}")

print(f"\nLearned weight: {regressor.coef_[0]:.4f}")
print(f"Learned bias:   {regressor.intercept_[0]:.4f}")
\end{codeholder}

\paragraph{Explication}
La régression est réalisée avec \texttt{SGDRegressor} (descente de gradient
stochastique). Les métriques rapportées incluent le MSE et le R\textsuperscript{2}.

\subsection*{Résultats et visualisation}

\begin{figure}[htbp]
    \centering
    \includegraphics[width=0.9\textwidth]{regression.png}
    \caption{Régression linéaire (jeu d'entraînement vs prédictions)
             et visualisation de la droite ajustée.}
    \label{fig:regression}
\end{figure}

\paragraph{Résultats chiffrés}
Lors de l'exécution :
\begin{itemize}
    \item Training MSE: 3.3933
    \item Test MSE: 2.5993
    \item Test MAE: 1.1658
    \item R\textsuperscript{2} (Test): 0.9826
    \item Poids appris: 11.3652 ; Biais appris: -0.3648
\end{itemize}

\paragraph{Interprétation}
Le modèle capture bien la relation linéaire malgré le bruit ajouté.
Le R\textsuperscript{2} proche de 0.98 indique un excellent ajustement sur
le jeu de test.


% ======================== EXERCICE 4 ========================
\exercice{Classification sur MNIST (digits)}

\subsection*{Énoncé}
Utilisez un jeu de données célèbre comme MNIST (option : utilisez un
sous-ensemble pour simplifier). Préparez les données pour un modèle Perceptron.
Entraînez le modèle sur un ensemble d'entraînement et testez-le sur un
ensemble de validation. Discutez des performances et des limitations du
Perceptron sur des données complexes.

\subsection*{1) Chargement et préparation}

\begin{codeholder}{ex4.py}
import numpy as np
import matplotlib.pyplot as plt
from sklearn.datasets import load_digits
from sklearn.model_selection import train_test_split
from sklearn.preprocessing import StandardScaler
from sklearn.linear_model import SGDClassifier
from sklearn.metrics import (
    accuracy_score, classification_report, confusion_matrix
)
import seaborn as sns

# Load MNIST digits dataset (subset: 0-9 digits, 1797 samples)
digits = load_digits()
X      = digits.data
y      = digits.target

print(f"Dataset shape:     {X.shape}")
print(f"Number of classes: {len(np.unique(y))}")
print(f"Classes:           {np.unique(y)}")

# Split data into training and test sets
X_train, X_test, y_train, y_test = train_test_split(
    X, y, test_size=0.2, random_state=42, stratify=y
)

# Standardize features
scaler         = StandardScaler()
X_train_scaled = scaler.fit_transform(X_train)
X_test_scaled  = scaler.transform(X_test)
\end{codeholder}

% --

\subsection*{2) Entraînement et évaluation}

\begin{codeholder}{ex4.py}
# Train Perceptron
clf = SGDClassifier(
    loss='perceptron',
    eta0=0.01,
    max_iter=100,
    random_state=42,
    tol=1e-3
)
clf.fit(X_train_scaled, y_train)

# Make predictions
y_pred_train = clf.predict(X_train_scaled)
y_pred_test  = clf.predict(X_test_scaled)

# Evaluate model
train_accuracy = accuracy_score(y_train, y_pred_train)
test_accuracy  = accuracy_score(y_test,  y_pred_test)

print(f"\nTraining Accuracy: {train_accuracy:.2%}")
print(f"Test Accuracy:     {test_accuracy:.2%}")

print("\nClassification Report (Test Data):")
print(classification_report(y_test, y_pred_test))
\end{codeholder}

\paragraph{Explication}
Le Perceptron (via \texttt{SGDClassifier} avec perte 'perceptron') est
appliqué aux vecteurs de pixels bruts. Les performances sont correctes
mais nettement inférieures à des réseaux convolutifs.

% --

\subsection*{4) Implémentation avec Keras}

\begin{codeholder}{ex4.py}
import numpy as np
import matplotlib.pyplot as plt
from sklearn.datasets import load_digits
from sklearn.model_selection import train_test_split
from sklearn.preprocessing import StandardScaler
from sklearn.metrics import accuracy_score, classification_report, confusion_matrix
import seaborn as sns
import tensorflow as tf
from tensorflow.keras.models import Sequential
from tensorflow.keras.layers import Dense
from tensorflow.keras.optimizers import SGD
from tensorflow.keras.utils import to_categorical

# Load MNIST digits dataset (subset: 0-9 digits, 1797 samples)
digits = load_digits()
X = digits.data
y = digits.target

print(f"Dataset shape:     {X.shape}")
print(f"Number of classes: {len(np.unique(y))}")
print(f"Classes:           {np.unique(y)}")

# Split data into training and test sets
X_train, X_test, y_train, y_test = train_test_split(
    X, y, test_size=0.2, random_state=42, stratify=y
)

# Standardize features
scaler = StandardScaler()
X_train_scaled = scaler.fit_transform(X_train)
X_test_scaled = scaler.transform(X_test)

# Convert labels to categorical for Keras (one-hot encoding)
y_train_cat = to_categorical(y_train, num_classes=10)
y_test_cat = to_categorical(y_test, num_classes=10)

# Create and train Perceptron using Keras (multi-class classification)
def create_perceptron_model(input_dim, num_classes):
    model = Sequential([
        Dense(num_classes, input_dim=input_dim, activation='softmax',  # Softmax for multi-class
              kernel_initializer='zeros', bias_initializer='zeros')
    ])
    optimizer = SGD(learning_rate=0.01)
    model.compile(optimizer=optimizer, loss='categorical_crossentropy', metrics=['accuracy'])
    return model

# Train the model
model = create_perceptron_model(X_train_scaled.shape[1], 10)
history = model.fit(X_train_scaled, y_train_cat, 
                    epochs=100, 
                    verbose=0,
                    validation_split=0.1)

# Make predictions
y_pred_train_prob = model.predict(X_train_scaled)
y_pred_train = np.argmax(y_pred_train_prob, axis=1)
y_pred_test_prob = model.predict(X_test_scaled)
y_pred_test = np.argmax(y_pred_test_prob, axis=1)

# Evaluate model
train_accuracy = accuracy_score(y_train, y_pred_train)
test_accuracy = accuracy_score(y_test, y_pred_test)

print(f"\nTraining Accuracy: {train_accuracy:.2%}")
print(f"Test Accuracy:     {test_accuracy:.2%}")

print("\nClassification Report (Test Data):")
print(classification_report(y_test, y_pred_test))

# Confusion matrix
cm = confusion_matrix(y_test, y_pred_test)
plt.figure(figsize=(10, 8))
sns.heatmap(cm, annot=True, fmt='d', cmap='Blues')
plt.title('Confusion Matrix for the test set (MNIST)')
plt.ylabel('True label')
plt.xlabel('Predicted label')
plt.tight_layout()
plt.savefig('mnist_confusion_matrix.png', dpi=100, bbox_inches='tight')
plt.show()

# Visualize some predictions
def plot_predictions(images, true_labels, pred_labels, num_samples=10):
    plt.figure(figsize=(12, 4))
    for i in range(num_samples):
        plt.subplot(2, num_samples//2, i+1)
        plt.imshow(images[i].reshape(8, 8), cmap='gray')
        color = 'green' if true_labels[i] == pred_labels[i] else 'red'
        plt.title(f'T:{true_labels[i]}\\nP:{pred_labels[i]}', color=color)
        plt.axis('off')
    plt.tight_layout()
    plt.savefig('mnist_predictions.png', dpi=100, bbox_inches='tight')
    plt.show()

# Show some test predictions
plot_predictions(X_test_scaled, y_test, y_pred_test, num_samples=10)

# Get weights and bias (for inspection)
weights = model.get_weights()[0]  # Shape: (64, 10) - 64 input features, 10 classes
bias = model.get_weights()[1]     # Shape: (10,)
print(f"\nWeights shape: {weights.shape}")
print(f"Bias shape: {bias.shape}")
\end{codeholder}

\paragraph{Explication}
Cette implémentation utilise Keras pour créer un perceptron simple avec une activation softmax pour la classification multi-classes.
Nous utilisons l'optimiseur SGD avec un taux d'apprentissage de 0.01 et optimisons la perte catégorielle croisée.
Les résultats montrent que l'approche Keras donne des performances similaires à celles de SGDClassifier de scikit-learn pour ce problème de classification.

% --

\subsection*{Résultats et visualisation}

\begin{figure}[htbp]
    \centering
    \includegraphics[width=0.8\textwidth]{mnist_confusion_matrix.png}
    \caption{Matrice de confusion pour le jeu de test (MNIST).}
    \label{fig:mnist_cm}
\end{figure}

\begin{figure}[htbp]
    \centering
    \includegraphics[width=0.8\textwidth]{mnist_predictions.png}
    \caption{Exemples d'images du jeu de test avec prédictions
             (vert = correct, rouge = incorrect).}
    \label{fig:mnist_pred}
\end{figure}

\paragraph{Résultats chiffrés}
Lors de l'exécution : \textbf{Training Accuracy = 99.03\%},
\textbf{Test Accuracy = 94.44\%}.

\paragraph{Interprétation}
Le Perceptron linéaire donne des résultats raisonnables sur MNIST
mais ne capture pas la structure spatiale des pixels ni les relations
non-linéaires complexes. Pour des performances d'état de l'art,
on utilisera des CNN.


% ======================== EXERCICE 5 ========================
\exercice{Segmentation des données synthétiques (make\_blobs)}

\subsection*{Énoncé}
Génération de données bidimensionnelles avec 3 centres, prétraitement,
implémentation d'un Perceptron simple (from-scratch) pour séparer deux
clusters, puis segmentation avec KMeans pour 3 clusters
et comparaison des résultats.

\subsection*{1) Génération et visualisation des données}

\begin{codeholder}{ex5.py}
import numpy as np
import matplotlib.pyplot as plt
from sklearn.datasets import make_blobs
from sklearn.model_selection import train_test_split
from sklearn.preprocessing import StandardScaler
from sklearn.cluster import KMeans

X, y = make_blobs(
    n_samples=300,
    centers=3,
    n_features=2,
    random_state=42,
    cluster_std=1.5
)

plt.figure(figsize=(10, 6))
plt.scatter(X[:, 0], X[:, 1], c=y, cmap='viridis', edgecolors='k', s=50)
plt.xlabel('Feature 1')
plt.ylabel('Feature 2')
plt.title('Generated Data with 3 Clusters (make_blobs)')
plt.colorbar(label='Cluster')
plt.grid(True, alpha=0.3)
plt.savefig('ex5_initial_data.png', dpi=100, bbox_inches='tight')
plt.show()

scaler   = StandardScaler()
X_scaled = scaler.fit_transform(X)

X_train, X_test, y_train, y_test = train_test_split(
    X_scaled, y, test_size=0.2, random_state=42, stratify=y
)
\end{codeholder}

% --

\subsection*{2) Implémentation du Perceptron (from scratch)}

\begin{codeholder}{ex5.py}
class Perceptron:

    def __init__(self, learning_rate=0.01, n_iterations=1000):
        self.learning_rate = learning_rate
        self.n_iterations  = n_iterations
        self.weights       = None
        self.bias          = None

    def fit(self, X, y):
        n_samples, n_features = X.shape
        self.weights = np.zeros(n_features)
        self.bias    = 0

        for _ in range(self.n_iterations):
            for idx, x_i in enumerate(X):
                linear_output = np.dot(x_i, self.weights) + self.bias
                y_predicted   = self._activation(linear_output)
                update        = self.learning_rate * (y[idx] - y_predicted)
                self.weights += update * x_i
                self.bias    += update

    def _activation(self, x):
        return np.where(x >= 0, 1, 0)

    def predict(self, X):
        linear_output = np.dot(X, self.weights) + self.bias
        return self._activation(linear_output)
\end{codeholder}

\paragraph{Explication}
La classe Perceptron implémente la règle d'apprentissage pour un modèle
linéaire binaire. On initialise les poids à zéro, puis pour chaque itération
on met à jour les poids selon l'erreur observée.

% --

\subsection*{2bis) Implémentation du Perceptron avec Keras}

\begin{codeholder}{ex5.py}
# Create and train Perceptron using Keras
def create_perceptron_model(input_dim):
    model = Sequential([
        Dense(1, input_dim=input_dim, activation='sigmoid',  # Sigmoid for binary classification
              kernel_initializer='zeros', bias_initializer='zeros')
    ])
    optimizer = SGD(learning_rate=0.1)
    model.compile(optimizer=optimizer, loss='binary_crossentropy', metrics=['accuracy'])
    return model

# Train the model
model = create_perceptron_model(X_train_binary.shape[1])
history = model.fit(X_train_binary, y_train_binary, 
                    epochs=1000, 
                    verbose=0,
                    validation_split=0.1)

# Make predictions
y_pred_train_prob = model.predict(X_train_binary)
y_pred_train = (y_pred_train_prob > 0.5).astype(int).flatten()
y_pred_test_prob = model.predict(X_test_binary)
y_pred_test = (y_pred_test_prob > 0.5).astype(int).flatten()

# Evaluate model
train_accuracy_keras = np.mean(y_pred_train == y_train_binary)
test_accuracy_keras = np.mean(y_pred_test == y_test_binary)

print(f"Training Accuracy (Keras): {train_accuracy_keras:.2%}")
print(f"Test Accuracy (Keras): {test_accuracy_keras:.2%}")

# Get weights and bias (for comparison with sklearn)
weights_keras = model.get_weights()[0].flatten()  # weights for the two features
bias_keras = model.get_weights()[1][0]           # bias
print(f"\nLearned weights (Keras): {weights_keras}")
print(f"Learned bias (Keras): {bias_keras}")

# Visualization
plt.figure(figsize=(10, 6))
h = 0.02
x_min, x_max = X_train[:, 0].min() - 1, X_train[:, 0].max() + 1
y_min, y_max = X_train[:, 1].min() - 1, X_train[:, 1].max() + 1
xx, yy = np.meshgrid(np.arange(x_min, x_max, h), np.arange(y_min, y_max, h))

Z = model.predict(np.c_[xx.ravel(), yy.ravel()])
Z = (Z > 0.5).astype(int).reshape(xx.shape)

plt.contourf(xx, yy, Z, alpha=0.3, cmap='RdBu')
plt.scatter(X_train_binary[y_train_binary == 0, 0], 
            X_train_binary[y_train_binary == 0, 1],
            c='blue', marker='o', label='Cluster 0', edgecolors='k')
plt.scatter(X_train_binary[y_train_binary == 1, 0], 
            X_train_binary[y_train_binary == 1, 1],
            c='red', marker='s', label='Cluster 1', edgecolors='k')
plt.xlabel('Feature 1 (Standardized)')
plt.ylabel('Feature 2 (Standardized)')
plt.title('Perceptron Decision Boundary (Keras) - Clusters 0 vs 1')
plt.legend()
plt.grid(True, alpha=0.3)
plt.savefig('ex5_perceptron_boundary_keras.png', dpi=100, bbox_inches='tight')
plt.show()
\end{codeholder}

\paragraph{Explication}
Nous définons ensuite notre modèle perceptron en utilisant Keras avec une activation sigmoïde adaptée à la classification binaire.
Nous utilisons l'optimiseur SGD avec un taux d'apprentissage de 0.1 et optimisons la perte binaire croisée.
Après avoir créé le modèle, nous l'entraînons sur les données d'entraînement standardisées.
Nous faisons ensuite des prédictions sur les ensembles d'entraînement et de test, puis nous évaluons le modèle en utilisant les mêmes métriques que précédemment : précision.
Nous extrayons également les poids et le biais appris pour les comparer avec ceux obtenus par l'implémentation from scratch.

% --

\subsection*{3) Segmentation non supervisée (KMeans) et comparaison}

\begin{codeholder}{ex5.py}
kmeans   = KMeans(n_clusters=3, random_state=42, n_init=10)
kmeans.fit(X_scaled)
y_kmeans = kmeans.labels_

plt.figure(figsize=(10, 6))
plt.scatter(
    X_scaled[:, 0], X_scaled[:, 1],
    c=y_kmeans, cmap='viridis', edgecolors='k', s=50, marker='o'
)
centers = kmeans.cluster_centers_
plt.scatter(
    centers[:, 0], centers[:, 1],
    c='red', marker='X', s=200, edgecolors='k', label='Centroids'
)
plt.xlabel('Feature 1 (Standardized)')
plt.ylabel('Feature 2 (Standardized)')
plt.title('KMeans Clustering (3 Clusters)')
plt.legend()
plt.grid(True, alpha=0.3)
plt.savefig('ex5_kmeans_clusters.png', dpi=100, bbox_inches='tight')
plt.show()

print("\n--- KMeans Results ---")
print(f"Inertia:              {kmeans.inertia_:.2f}")
print(f"Number of iterations: {kmeans.n_iter_}")
\end{codeholder}

\paragraph{Explication}
KMeans est utilisé pour obtenir une segmentation non supervisée en 3 clusters.
Les résultats sont comparés visuellement et via métriques comme l'inertie.

\subsection*{Résultats et visualisation}

\begin{figure}[htbp]
    \centering
    \includegraphics[width=0.6\textwidth]{ex5_initial_data.png}
    \caption{Données générées avec \texttt{make\_blobs} (3 clusters).}
    \label{fig:ex5_initial}
\end{figure}

\begin{figure}[htbp]
    \centering
    \includegraphics[width=0.45\textwidth]{ex5_perceptron_boundary.png}
    \hfill
    \includegraphics[width=0.45\textwidth]{ex5_perceptron_boundary_keras.png}
    \caption{Gauche : frontière de décision du Perceptron from scratch (clusters 0 vs 1).
             Droite : frontière de décision du Perceptron Keras (clusters 0 vs 1).}
    \label{fig:ex5_compare_perceptrons}
\end{figure}

\begin{figure}[htbp]
    \centering
    \includegraphics[width=0.45\textwidth]{ex5_perceptron_boundary.png}
    \hfill
    \includegraphics[width=0.45\textwidth]{ex5_kmeans_clusters.png}
    \caption{Gauche : frontière de décision du Perceptron (from scratch, clusters 0 vs 1).
             Droite : clusters trouvés par KMeans (n=3).}
    \label{fig:ex5_compare}
\end{figure}

\begin{figure}[htbp]
    \centering
    \includegraphics[width=0.9\textwidth]{ex5_comparison.png}
    \caption{Comparaison visuelle : données originales / Perceptron (binaire)
             / KMeans (3 clusters).}
    \label{fig:ex5_comp_full}
\end{figure}

\paragraph{Résultats chiffrés}
Lors de l'exécution :
\begin{itemize}
    \item Perceptron from scratch (clusters 0 vs 1) :
          Training Accuracy = 100.00\%, Test Accuracy = 100.00\%.
    \item Perceptron Keras (clusters 0 vs 1) :
          Training Accuracy = 100.00\%, Test Accuracy = 100.00\%.
    \item KMeans : Inertia = 39.07 ; nombre d'itérations = 2.
\end{itemize}

\paragraph{Interprétation}
Le Perceptron binaire sépare parfaitement les deux premiers clusters choisis.
KMeans retrouve 3 groupes proches de la vérité terrain ;
l'inertie donne une idée de la compacité des clusters.


% ======================== ANNEXES ========================
\clearpage
\section*{Annexes}
\addcontentsline{toc}{section}{Annexes}

\subsection*{Fichiers source complets}
Dans cette annexe, vous trouverez la liste des fichiers Python fournis
pour le TP. Pour garantir la traçabilité, les fichiers complets sont
joints au dossier de remise.

\vspace{1cm}
\textbf{Remarques finales}
\begin{itemize}
    \item Les extraits de code ci-dessus sont copiés strictement tels
          qu'ils apparaissent dans les fichiers fournis.
          Aucune modification du code des exercices n'a été effectuée.
    \item Les scripts Python peuvent être exécutés directement pour
          reproduire les figures et les résultats numériques présentés
          dans ce rapport.
\end{itemize}

\end{document}
